\section{Actividad 4. Análisis comparativo}
\subsection{Tecnologías de virtualización}
Las cifras siguientes son rangos orientativos: dependen del hardware, la imagen y la configuración. La madurez es una valoración de 1 a 5, no una medición estandarizada.
\begin{longtable}{p{2.2cm}p{2cm}p{2cm}p{2cm}p{2cm}p{2cm}}
\textbf{Criterio} & \textbf{VM} & \textbf{Contenedor} & \textbf{MicroVM} & \textbf{Wasm} & \textbf{Unikernel} \\
\hline
Arranque & Segundos a minutos & Milisegundos a segundos & Decenas a cientos de ms & Milisegundos o menos & Decenas de ms a segundos \\
Tamaño & GB & MB a GB & Decenas de MB o más & KB a MB, sin contar runtime & MB a decenas de MB \\
Aislamiento & Kernel propio e hipervisor & Procesos y namespaces; kernel compartido & Kernel propio y VMM reducido & Sandbox del runtime; capacidades limitadas & Imagen especializada, normalmente aislada por hipervisor \\
Ejemplo & KVM/Proxmox VE & Docker & Firecracker & Wasmtime & Unikraft \\
Uso & Sistemas heredados y SO distintos & Servicios web y despliegue continuo & Cargas efímeras multiinquilino & Extensiones y edge & Servicios especializados \\
Madurez & 5 & 5 & 4 & 3 & 2 \\
\end{longtable}
Las categorías no son excluyentes: un contenedor puede correr dentro de una VM y un módulo Wasm dentro de un contenedor. Un unikernel no garantiza aislamiento por sí solo: depende del entorno en que se ejecute. Los rangos son comparativos, no garantías de rendimiento.

\subsection{Broadcom y VMware}
\begin{itemize}
\item \textbf{Compra.} Broadcom anunció el cierre el 22 de noviembre de 2023; el valor citado de la operación fue aproximadamente 69 mil millones de dólares. El valor anunciado de una operación en efectivo y acciones puede diferir del valor final según la fecha y el precio de las acciones. \href{https://investors.broadcom.com/news-releases/news-release-details/broadcom-completes-acquisition-vmware}{Broadcom, 2023}.
\item \textbf{Incrementos.} La práctica menciona 150\%, pero no se encontró una muestra verificable que permita generalizar esa cifra como promedio mundial. Hay casos muy distintos según contratos y productos; por ejemplo, en Países Bajos se reportó un incremento de 85\% para Rijkswaterstaat. Conviene comparar propuestas concretas, no extrapolar un caso. \href{https://www.networkworld.com/article/4015489/dutch-court-forces-broadcom-to-support-vmware-migration-after-85-price-hike-backlash.html}{Network World, 2025}.
\item \textbf{Pruebas 56 de 57.} La alternativa fue Proxmox VE. El estudio de Blockbridge comparó configuraciones de almacenamiento NVMe/TCP bajo cargas específicas; el resultado no demuestra que Proxmox sea superior en todo sistema ni para toda carga. \href{https://kb.blockbridge.com/technote/proxmox-vs-vmware-nvmetcp/}{Blockbridge}.
\item \textbf{Pronóstico.} Una predicción de Gartner difundida en 2025 estimó que 35\% de las cargas VMware migrarían a otras plataformas para 2028; es un pronóstico de cargas, no una medición de clientes que ya migraron. \href{https://arstechnica.com/information-technology/2025/09/35-percent-of-vmware-workloads-expected-to-migrate-elsewhere-by-2028/}{Ars Technica, 2025}.
\item \textbf{Litigio.} En Países Bajos, la agencia pública Rijkswaterstaat obtuvo una resolución para mantener soporte durante una migración. El caso se relaciona con condiciones de soporte y continuidad, no demuestra por sí solo una sanción general por precios excesivos. \href{https://www.networkworld.com/article/4015489/dutch-court-forces-broadcom-to-support-vmware-migration-after-85-price-hike-backlash.html}{Network World, 2025}.
\item \textbf{Decisión para 50 VM.} Inventariar contratos, dependencias, respaldos y requisitos de disponibilidad; comparar el costo total de renovar con Proxmox, Hyper-V y nube durante tres a cinco años. Haría un piloto con cargas no críticas, probaría restauraciones y rendimiento y migraría por etapas con reversión. Mantendría temporalmente en VMware los sistemas cuya compatibilidad o costo de migración no justifique cambiarlos.
\end{itemize}

\subsection{Preguntas de análisis}
\begin{enumerate}
\item \textbf{Kubernetes y Swarm.} Kubernetes ofrece reconciliación declarativa, actualizaciones graduales con rollback y extensibilidad mediante recursos y controladores; también dispone de capacidades amplias para políticas, almacenamiento y redes. En el ecosistema influyen el respaldo de CNCF y los proveedores de nube, además de herramientas, documentación y personal capacitado. Swarm conserva valor por su sencillez, pero el ecosistema de Kubernetes es más amplio.
\item \textbf{Wasm y Docker.} La afirmación de Hykes destaca la portabilidad de módulos y su aislamiento en un runtime compacto: parte del problema de distribuir una aplicación podría abordarse de otra forma. Sin embargo, Wasm no sustituye automáticamente dependencias nativas, redes, persistencia ni todas las aplicaciones existentes. Es probable que complemente a los contenedores en edge, extensiones y funciones breves, mientras estos siguen siendo útiles para software Linux convencional.
\item \textbf{Confidential computing.} Protege datos durante su procesamiento mediante entornos de ejecución confiables respaldados por hardware. Tecnologías como AMD SEV-SNP aíslan y cifran memoria de VM; Azure Confidential VMs ofrece opciones basadas en estas tecnologías. En IA puede ayudar a proteger datos y modelos sensibles durante inferencia o entrenamiento, aunque no elimina riesgos de configuración, aplicación o datos de salida. \href{https://learn.microsoft.com/en-us/azure/confidential-computing/confidential-vm-overview}{Microsoft, Confidential VMs}.
\end{enumerate}

\section{Actividad 5. El futuro de la virtualización: de las máquinas virtuales a los microsegundos}

La virtualización ha cambiado de forma conforme cambian las necesidades de la industria. Su historia no es una carrera en la que cada tecnología elimina a la anterior, sino una sucesión de maneras de repartir recursos y aislar cargas. En los mainframes se buscaba aprovechar equipos costosos entre varios usuarios; las máquinas virtuales extendieron esa idea a servidores comunes, con sistemas operativos independientes sobre un hipervisor. Después, los contenedores permitieron empaquetar procesos y dependencias compartiendo un kernel. Las microVM redujeron el costo de crear huéspedes aislados y WebAssembly incorporó otro nivel: módulos portables ejecutados por un runtime con acceso controlado a recursos. El patrón es una abstracción cada vez más cercana a la aplicación, aunque cada capa conserva un papel propio.

Una VM sigue siendo adecuada cuando se necesita un kernel independiente, compatibilidad con un sistema operativo concreto o una frontera de aislamiento fuerte. El precio es administrar imágenes y sistemas invitados, además de reservar memoria y almacenamiento. Los contenedores suelen iniciar más rápido y facilitan que desarrollo y producción ejecuten paquetes semejantes. Su portabilidad no es absoluta: dependen de la arquitectura, del kernel y de servicios externos. En una Mac, por ejemplo, los contenedores Linux necesitan una capa Linux, normalmente proporcionada por una VM. Comprender esa dependencia evita presentar a Docker y a la virtualización de hardware como opciones incompatibles.

La orquestación convirtió a los contenedores en una plataforma para aplicaciones distribuidas. Kubernetes mantiene el estado deseado, reemplaza pods, escala réplicas y automatiza actualizaciones. La encuesta CNCF publicada en 2026 reportó que 82\% de los usuarios de contenedores encuestados ejecutaba Kubernetes en producción. La cifra muestra su adopción dentro de esa población, no de todas las empresas del mundo. Su éxito también acarrea costos: redes, observabilidad, seguridad y actualizaciones requieren personal y tiempo. Para una aplicación pequeña, un servicio gestionado o incluso una VM bien administrada pueden ser más simples.

Las microVM buscan combinar tiempos de arranque reducidos con un kernel separado. Firecracker, por ejemplo, está diseñado para ejecutar cargas de corta duración con una superficie de dispositivo reducida. Wasm presenta una propuesta distinta: distribuir módulos compactos y limitar las capacidades que reciben del entorno. Puede ser útil para extensiones, procesamiento cercano al usuario y funciones breves, pero sus interfaces de sistema, bibliotecas y herramientas todavía condicionan los casos de uso. Un módulo Wasm puede ejecutarse junto con contenedores; por eso no considero inevitable que los reemplace.

La compra de VMware por Broadcom, concluida en 2023, agregó un factor económico a la decisión técnica. La transición de licencias perpetuas a suscripciones y paquetes impulsó a clientes a revisar sus contratos y evaluar otras plataformas. No todos enfrentaron el mismo aumento y migrar implica costos de personal, pruebas, respaldo y posibles interrupciones. En consecuencia, evaluar Proxmox u otra alternativa abierta no equivale a migrar inmediatamente. Un piloto con cargas representativas y un plan de reversión ofrece evidencia más útil que asumir que una plataforma gana en todas las pruebas. Para una empresa con 50 VM, la prioridad debe ser conocer dependencias y medir el costo total antes de programar ventanas de cambio.

Durante los próximos cinco años espero que las aplicaciones web nuevas continúen utilizando contenedores y plataformas gestionadas, con Kubernetes donde la escala o la necesidad de control justifiquen su complejidad. Para IA y aprendizaje automático predominará una combinación: contenedores para distribuir software, orquestadores para programar trabajos y mecanismos específicos para compartir o asignar GPU. El aislamiento confidencial tendrá más valor cuando se procesen datos delicados, pero no solucionará por sí mismo filtraciones causadas por una aplicación mal diseñada. Para IoT y edge, Wasm puede ganar espacio por su tamaño y ejecución controlada, mientras Linux y contenedores seguirán presentes en dispositivos más capaces. Las aplicaciones heredadas continuarán usando VM mientras su sistema operativo y sus proveedores requieran ese entorno.

En México estas decisiones deben relacionarse con conectividad, regulación, costos y personal disponible. La región AWS México Central, identificada como \texttt{mx-central-1}, permite considerar infraestructura local frente a regiones de otros países. Su presencia no garantiza automáticamente menor costo ni cumplimiento normativo: hay que medir latencia, revisar disponibilidad de servicios y diseñar respaldos y permisos. Una organización mexicana mediana puede beneficiarse de servicios gestionados sin comprometer todo su sistema a un único proveedor si conserva procesos de exportación y recuperación probados.

Para quienes egresan del IPN, la ventaja profesional no consiste solamente en conocer comandos. Es necesario entender redes, Linux, programación, monitoreo y seguridad, así como explicar por qué se eligió cada tecnología. La práctica de probar recuperación, comparar imágenes y registrar resultados desarrolla ese criterio. También enseña a distinguir una medida de laboratorio de una promesa universal: el tiempo de recuperación de un pod en una computadora no garantiza el comportamiento de un sistema de producción. El futuro será híbrido: distintas abstracciones convivirán y el trabajo de ingeniería consistirá en elegir la combinación que resuelva el problema con un costo y un riesgo aceptables.

\subsection{Fuentes para el análisis y el ensayo}
\begin{itemize}
\item CNCF. (2026). \href{https://www.cncf.io/wp-content/uploads/2026/01/CNCF_Annual_Survey_Report_final.pdf}{Annual Cloud Native Survey 2025}.
\item AWS. (2025). \href{https://docs.aws.amazon.com/global-infrastructure/latest/regions/doc-history.html}{Regions and Availability Zones: document history}.
\item AWS. (2026). \href{https://docs.aws.amazon.com/pdfs/whitepapers/latest/security-design-of-aws-nitro-system/security-design-of-aws-nitro-system.pdf}{The Security Design of the AWS Nitro System}.
\item Microsoft. (2026). \href{https://learn.microsoft.com/en-us/azure/confidential-computing/confidential-vm-overview}{Confidential virtual machines}.
\item Ars Technica. (2025). \href{https://arstechnica.com/information-technology/2025/09/35-percent-of-vmware-workloads-expected-to-migrate-elsewhere-by-2028/}{Gartner y migración de cargas VMware}.
\item Network World. (2025). \href{https://www.networkworld.com/article/4015489/dutch-court-forces-broadcom-to-support-vmware-migration-after-85-price-hike-backlash.html}{Resolución sobre soporte para Rijkswaterstaat}.
\item Broadcom. (2023). \href{https://investors.broadcom.com/news-releases/news-release-details/broadcom-completes-acquisition-vmware}{Cierre de la adquisición de VMware}.
\end{itemize}
